\section*{v4.7.0 Full-Monolith Release Scope and Provenance Notice}
\addcontentsline{toc}{section}{v4.7.0 Full-Monolith Release Scope and Provenance Notice}
\label{sec:v470-rc9-full-monolith-recovery-notice}

\noindent\textbf{Recovery status.} This v4.7.0 full-monolith source preserves the complete working monolith by using the last full-length source layer r54 as the base, appending the post-r54 guarded solve chain through r80, and retaining the r81--r84 frame-typed effective/local closure clarifications.  Earlier compact upload handoff artefacts are not release bases because they compiled only a short consolidation artefact rather than the full monolith.  The present release-candidate source keeps the corrected full-source recovery build and the compact title-page discipline.

\noindent\textbf{DOI status.} The intended reserved v4.7.0 Version DOI is \href{https://doi.org/10.5281/zenodo.20232911}{10.5281/zenodo.20232911}. The Companion Concept DOI is \href{https://doi.org/10.5281/zenodo.19210728}{10.5281/zenodo.19210728}. The previous/baseline v4.6.0 DOI remains \href{https://doi.org/10.5281/zenodo.20117361}{10.5281/zenodo.20117361} as provenance only.

\noindent\textbf{Claim-scope lock.} The release scope is guarded local/restricted only. It allows the local/restricted downstream chain, r60+ conservative C4 relative-seam neutralization, C5 detector-threshold terminality at \(\theta=1/256\) after \(n=20\), local photon-clock selector preservation under \(\Delta_\gamma<1/5\), typed mixed-comparison local control, terminal-data bridge, restricted terminal-data naturality, and terminal blocker freeze. It does not claim exact finite C5 zero closure, exact all-readout terminality, C6/T2 closure, a global photon-clock/RH theorem, or global all-readout closure.






\section*{v4.7.0-r44 Internal Working Status}
\addcontentsline{toc}{section}{v4.7.0-r44 Internal Working Status}
\noindent\textbf{Branch status.} This source is an internal source-only working
candidate over the v4.6.0 release-final baseline and the v4.7.0-r43 local
certificate handoff packet. It performs a source-only handoff integrity scan
and builds a local certificate index for the next working layer. It assigns no new Zenodo version DOI, authorizes no public
upload, and produces no PDF in this intermediate pass.

\noindent\textbf{Mathematical status.} The r44 pass is a handoff-integrity and
local-certificate-index pass. It does not promote the local MM/HC monitor certificate
to a physical action, a Lagrangian, a derived field equation, a global readout
closure, or a T2 theorem. It keeps the r43 guard vector active: \(C4\) remains
route-stable with \(\beta_4=1/256\), finite \(C5\)-closure is not claimed,
\(C6/T2\) remains blocked, the RH channel remains a photon-clock handoff probe
only, and global all-readout closure remains unclaimed.

\noindent\textbf{Reader rule.} Everything below should be read as a local
MM/HC interface certificate extracted from a guarded atlas handoff.  The
resulting objects are certified local interface constraints, not derived
physical laws.
\clearpage
\section*{v4.6.0 Release-Final Baseline Status}
\addcontentsline{toc}{section}{v4.6.0 Release-Final Baseline Status}
\noindent\textbf{Release-final baseline note.} This inherited baseline block records v4.6.0 release final with the version DOI \href{https://doi.org/10.5281/zenodo.20117361}{10.5281/zenodo.20117361}. The main title remains \emph{Quantum Sphaera Companion}. In the public v4.6.0 release-final source, the title page kept the Concept DOI and the Version DOI as separate DOI records. In this v4.7.0-r44 internal working source, the same DOI is listed as the baseline DOI and no new DOI is assigned. No all-readout/T2 closure is claimed.

\noindent\textbf{Main title rule.} The main title remains \emph{Quantum Sphaera Companion}.  The phrase \emph{Bounded Readout Closure and T2 Scope Boundary} is a release-theme subtitle only, not a replacement for the main title.

\medskip
\noindent\textbf{Baseline DOI-line rule.} The v4.7.0-r44 title page carries the stable concept DOI and exactly one baseline version DOI line: \emph{Baseline Version DOI (v4.6.0): 10.5281/zenodo.20117361}.  The concept DOI remains the stable Companion-line pointer.  The date line records the internal working-candidate status, not a new public DOI assignment.

\medskip
\noindent\textbf{Release-final scope.} This integrated release-final source records the v4.6.0 bounded readout result and the T2 scope boundary under the assigned Version DOI \href{https://doi.org/10.5281/zenodo.20117361}{10.5281/zenodo.20117361}.  It is the public release-final metadata state for v4.6.0, but it is not an all-readout closure theorem and does not promote any claim outside the declared bounded grammar.

\medskip
\noindent\textbf{Governing v4.6.0 status.}
\[
\mathsf{BRead}_{\mathrm{bounded}}
\in
\mathsf{Cons}(\mathcal G_{\mathrm{global}}^{\mathrm{decl}}),
\qquad
\mathsf{BRead}_{\mathrm{bounded}}\not\Rightarrow T2,
\qquad
T2\in\mathsf{Ext}(\mathcal G_{\mathrm{global}}^{\mathrm{decl}}).
\]
No empirical validation, Standard-Model-spectrum derivation, unrestricted physical readout, or all-readout closure is claimed.



\clearpage
\section*{Human-Facing Reading Layer and Notation Discipline}
\addcontentsline{toc}{section}{Human-Facing Reading Layer and Notation Discipline}

This edition separates the mathematical object language from revision and audit metadata.
Internal revision identifiers and combined version--revision
strings are not part of the physics or mathematics.  They are provenance tags only.
They may occur in file names, manifests, checksums, audit reports, and reconstruction
notes, but they are not used as exponents, object names, states, fields, obstructions,
or physical regimes in the reader-facing text.

The main reading order is therefore:
\begin{enumerate}[label=\arabic*.]
    \item the structural and physical claim in ordinary mathematical notation;
    \item the local hypotheses, guards, and return conditions;
    \item the obstruction or witness that remains after the permitted tests;
    \item the machine-readable audit trail, kept only as provenance material.
\end{enumerate}

When a previous draft used a revision-bearing symbol, this edition rewrites
it as a semantic object such as \(X^{\mathrm{rtg}}\), \(X^{\mathrm{lock}}\),
\(X^{\mathrm{nonsplit}}\), or simply \(X\), according to the mathematical role of the
object.  Version numbers remain valid release metadata, not mathematical structure.

\medskip
\noindent\textbf{Reader rule.}
The human-facing document should be readable without knowing the internal patch history.
The audit/provenance layer is still preserved for reconstruction and LLM-assisted
maintenance, but it is subordinate to the mathematical presentation.



\clearpage
\section*{Reader-Facing Proof Spine and Document Routing}
\addcontentsline{toc}{section}{Reader-Facing Proof Spine and Document Routing}

This v4.6.0 release-final frontmatter treats the mathematical narrative as the primary document and the machine/audit trail as secondary provenance.  The inherited v4.5.0 release-final layer remains visible as prior release genealogy, but the governing public metadata state is now v4.6.0 with Version DOI \href{https://doi.org/10.5281/zenodo.20117361}{10.5281/zenodo.20117361}.  The reader should not need the patch history in order to understand the current object language.  Revision identifiers, local build labels, and version-transition notes are therefore moved out of the proof spine and into the provenance layer.

\subsection*{Primary reading order}
\addcontentsline{toc}{subsection}{Primary reading order}

\begin{enumerate}[label=\arabic*.]
    \item \textbf{Structural claim.}  Read the HC/MM/chart interface as a structural bridge problem with explicit guards and comparison norms.
    \item \textbf{Local grammar.}  The photon-frame HC-halo/Dixon/lens grammar gives the inherited local admissible test surface.
    \item \textbf{Finite tests.}  Candidate menus are used only as certified finite screens.  They may consume, route, block, or name an obstruction; they do not create a physical claim by themselves.
    \item \textbf{Current v4.6.0 endpoint.}  The active endpoint is the bounded-readout surface: it is conservative and debt-free inside the declared grammar, but it does not imply T2 and does not discharge all possible readouts.
    \item \textbf{Return and handoff contract.}  A future all-readout answer must return one of three objects: a kill certificate, a strict norm-decrease witness, or a named persistent obstruction.  The v4.7 handoff therefore starts with the MM/HC exterior--mixed field-equation program, not with a T2 closure claim.
\end{enumerate}

\subsection*{Current semantic object registry}
\addcontentsline{toc}{subsection}{Current semantic object registry}

The active state is represented by semantic names, not revision names.  The older route-to-global bridge symbols remain valid v4.5 genealogy; the v4.6.0 governing objects are the bounded-readout surface, the all-readout scope boundary, and the MM/HC exterior--mixed handoff:
\begin{center}
\small
\begin{tabular}{@{}>{$}l<{$} P{0.25\textwidth} P{0.45\textwidth}@{}}
\toprule
\text{Symbol} & \textbf{Role} & \textbf{Reader meaning}\\
\midrule
B_{\mathrm{tie}} & inherited chart-tie component & v4.5 non-split coupling extracted from the local pair\\
B_{\mathrm{lock}} & inherited locked remnant & v4.5 remaining tie object after partial cancellation screens\\
\Obs_{\Bridge}^{\mathrm{lock}} & inherited locked bridge obstruction & v4.5 bridge obstruction; genealogy, not the final v4.6 endpoint\\
\mathsf{BRead}_{\mathrm{bounded}} & bounded readout surface & declared bounded-readout object: conservative inside the declared grammar\\
\Theta^{\mathrm{all}} & all-readout demand & unrestricted readout family required for T2; not supplied by the bounded result\\
\mathfrak O_{\mathrm{all}}^{4.6.0} & frozen all-readout obstruction & missing universal exterior selector, all-readout naturality family, and unrestricted exterior collar slot\\
\mathcal M_{\mathrm{bd}},\mathcal M_{\mathrm{ext}},\mathcal M_{\mathrm{mix}} & Millennium-matrix split & bounded, exterior, and mixed parts for the next field-equation determination\\
\mathcal H_{\mathrm{bd}},\mathcal H_{\mathrm{ext}},\mathcal H_{\mathrm{mix}} & HC split & bounded, exterior, and mixed compensator portions matched to the same readout split\\
\mathfrak E_{\mathrm{decl}} & field-equation reset template & structure-first equation template for the v4.7 MM/HC exterior--mixed program\\
\bottomrule
\end{tabular}
\end{center}

\subsection*{Guard discipline}
\addcontentsline{toc}{subsection}{Guard discipline}

No reader-facing claim is promoted unless the relevant guard is explicitly discharged.  In particular, this edition does not claim T2, all-readout closure, global OP closure, a global \(HS^2\) focal-return theorem, a through-collapse trajectory, an unverified atlas expansion, or a Standard-Model-spectrum derivation.  The bounded-readout theorem is therefore read as a conservative declared-grammar result, while \(\Theta^{\mathrm{all}}\) remains a scope boundary unless a separate all-readout witness family is supplied.

\subsection*{Human/LLM layering rule}
\addcontentsline{toc}{subsection}{Human/LLM layering rule}

The main document contains the mathematical object language.  The appendix may contain provenance: source packets, diffs, checksums, machine-readable registers, build logs, and reconstruction notes.  Provenance may mention build labels.  The proof spine must not use them as mathematical exponents, object names, or physical states.


\subsection*{Compressed chronology reading rule}
\addcontentsline{toc}{subsection}{Compressed chronology reading rule}

For reader navigation, inherited working chronology is to be compressed into semantic milestones rather than read as active mathematical notation.  The compression convention is:
\begin{center}
\small
\begin{tabular}{@{}p{0.16\textwidth}p{0.46\textwidth}p{0.26\textwidth}@{}}
\toprule
chronology stratum & reader-facing meaning & status\\
\midrule
$\mathrm{C}_0$ & earlier local readout and release genealogy & provenance\\
$\mathrm{C}_1$ & v4.4 declared-local OP1/OP4 closure ledger & prior release result\\
$\mathrm{C}_2$ & v4.5.0 declared global-boundary synthesis & prior theorematic release layer\\
$\mathrm{C}_3$ & v4.5.0 DOI/title/frontmatter repair & release metadata\\
$\mathrm{C}_4$ & v4.5.x editorial QA and chronology compression & post-release readability layer\\
$\mathrm{C}_5$ & v4.6.0 bounded-readout theorem and T2 scope boundary & current governing release status\\
$\mathrm{C}_6$ & MM/HC exterior--mixed field-equation determination & v4.7 handoff program\\
\bottomrule
\end{tabular}
\end{center}
Old working strings such as release-candidate notes, prior version DOIs, and local revision labels are therefore genealogy unless a nearby release-final or governing-status block explicitly says otherwise.  This rule changes only reader navigation; it does not change labels, references, DOI records, theorem statements, or scope guards.

\paragraph{Proof-spine phrase audit rule.}
In running prose, build expressions such as source-native rebuild, release candidate, freeze audit, packet, and working layer are reader-facing provenance terms unless they are explicitly part of a displayed ledger.  Where such terms occur in the active proof spine, the preferred reader form names the mathematical role first---for example grammar exhaustion, readout invariance, debt-free boundary transport, no-new-debt extension, or functorial comparison stability---and keeps the build phrase only as provenance.  This rule does not rename labels, change theorem statements, alter DOI records, or change any scope guard.

\subsection*{Reader-facing phrase-normalization rule}
\addcontentsline{toc}{subsection}{Reader-facing phrase-normalization rule}

Current-line section headings are allowed to use semantic names even when the
source labels and audit ledgers retain their original revision identifiers.
Thus a heading such as a witness execution, boundary audit, source-native
checkpoint, or title-page repair is to be read by its mathematical function,
not by the working-string that produced it.  Revision strings remain valid in
labels, comments, diffs, ledgers, and reconstruction records, but they are not
mathematical object names, proof data, tensor indices, route states, or closure
witnesses.


\thispagestyle{fancy}


\begin{center}
\fbox{\begin{minipage}{0.92\textwidth}
\small
\textbf{Companion status.} This is the v4.6.0 release-final frontmatter after bounded-readout integration, T2 scope-boundary freeze, and DOI normalization.  It preserves the declared mathematical global-boundary closure audit path from the v4.5.0 release line as provenance while making the v4.6.0 Version DOI \href{https://doi.org/10.5281/zenodo.20117361}{10.5281/zenodo.20117361} the current public release DOI. Version labels are retained only as provenance metadata; the mathematical proof spine uses semantic names.

\medskip
\noindent\textbf{Governing theorematic status.} Inside the declared global-boundary synthesis grammar
\[
 \mathcal G_{\mathrm{global}}^{\mathrm{decl}},
\]
OP~1, OP~4, and OP~5 are closed as a declared mathematical global-boundary closure. This does not claim unrestricted all-readout closure, global physical closure, QFT/GR closure, QCD/confinement, CKM/PMNS, quark-mass, or Standard-Model-spectrum derivation.

\medskip
\noindent\textbf{DOI rule.} The concept DOI remains the stable title-page pointer. The v4.5.0 version-specific DOI is \href{https://doi.org/10.5281/zenodo.20079410}{10.5281/zenodo.20079410}; prior version-specific DOI records are provenance metadata only. Version IDs, DOI fragments, dates, and release IDs are not mathematical object names, indices, theorem labels, route tags, or proof data.

\medskip
\noindent\textbf{Internal r10 proofread status.} This v4.5.1-r10 layer is an internal selected-range proofread and metadata-hygiene pass. It does not authorize public upload, does not insert a new DOI, does not emit Zenodo release text, and does not change the declared mathematical closure scope.

\medskip
\noindent\textbf{Internal r11 terminology-sweep status.} This v4.5.1-r11 layer is an internal targeted terminology sweep over the older OP5 and OP1/OP4 compatibility pages. It does not authorize public upload, does not insert a new DOI, does not emit Zenodo release text, does not rename theorem labels, and does not change the declared mathematical closure scope.

\medskip
\noindent\textbf{Internal r13 glossary/index intake status.} This v4.5.1-r13 layer is an internal glossary and index intake for the declared global-boundary synthesis grammar and the OP5/OP1/OP4 terminology family.  It is a reader-routing layer only: no public upload is authorized, no new DOI is inserted, no Zenodo release text is emitted, no theorem label is renamed, and no mathematical closure scope is changed.

\medskip
\noindent\textbf{Internal r13 glossary-appendix worksheet status.} This v4.5.1-r13 layer is an internal extraction worksheet for a possible future glossary appendix. It does not create a public release refresh, does not insert a new DOI, does not emit Zenodo release text, does not rename theorem labels, and does not change the declared mathematical closure scope.
\end{minipage}}
\end{center}
\vspace{0.6em}

\begin{abstract}
This v4.6.0 release-final frontmatter records the Quantum Sphaera Companion after bounded-readout integration, T2 scope-boundary freeze, DOI synchronization, and title-page/frontmatter hygiene.  The version-specific DOI for this release line is \href{https://doi.org/10.5281/zenodo.20117361}{10.5281/zenodo.20117361}.  The Companion concept DOI remains \href{https://doi.org/10.5281/zenodo.19210728}{10.5281/zenodo.19210728}.  The v4.5.0 version DOI \href{https://doi.org/10.5281/zenodo.20079410}{10.5281/zenodo.20079410} is retained as prior-release provenance.

The governing v4.6.0 release-facing result is the bounded-readout theorem together with the T2 scope-boundary statement.  Inside the declared global-boundary grammar \(\mathcal G_{\mathrm{global}}^{\mathrm{decl}}\), the bounded readout surface is conservative and debt-free, while T2 remains an explicit all-readout extension problem rather than a theorem asserted by this release.

This is a scoped mathematical theorem.  It is not a global physical closure theorem, not a global QFT/GR theorem, not unrestricted closure over every possible readout or physical model, and not a T2 closure claim.  Earlier version-specific DOI records remain provenance metadata only and are not mathematical object names, theorem labels, route tags, or proof data.
\end{abstract}

\clearpage
\section*{Frontmatter reader architecture}
\addcontentsline{toc}{section}{Frontmatter reader architecture}
\begin{quote}\small
\noindent\textbf{Purpose.} This page separates the current release-facing route from historical theorem memory. It is an orientation layer only; it does not replace any definition, theorem, proof, audit ledger, or historical witness later in the document.
\end{quote}

\begin{table}[H]
\centering
\small
\begin{tabular}{>{\raggedright\arraybackslash}p{0.24\textwidth} >{\raggedright\arraybackslash}p{0.31\textwidth} >{\raggedright\arraybackslash}p{0.31\textwidth}}
\toprule
\textbf{Layer} & \textbf{Human reading route} & \textbf{LLM/audit route}\\
\midrule
Current reader status & Read this as the v4.6.0 release-final frontmatter after bounded-readout integration, T2 scope-boundary freeze, DOI synchronization, and metadata hygiene. & Treat this pass as reader-facing organization and provenance separation only; no new mathematics is introduced beyond the audited v4.6.0 bounded-readout layer.\\
OP5 status & Read OP~5 as locally closed only inside the declared photon-frame HC-halo/Dixon grammar. & Track the required zero registers \(L_{\mathrm{unr}},B_{\mathrm{HC/MM}},R_{\mathrm{HS4}},D_{\mathrm{HS5}}\) and reject any stronger global reading.\\
QAM/MML/MMA & Treat QAM as store/object grammar, MML as routing/command grammar, and MMA as execution/readout semantics. & Use finite object bridges, instruction alphabets, deterministic guarded transitions, and route-status tables as audit handles.\\
Dixon coding & Read the Dixon layer as a 64-dimensional carrier with a \(2^{64}\)-mask activation space. & Keep carrier dimension and Boolean mask count separate; do not read Dixon itself as \(2^{64}\)-dimensional.\\
Historical witness & Use it as genealogy and reconstruction support. & Preserve it as non-destructive proof memory; do not read inherited baseline language as the current public claim unless explicitly mapped by the local proof chain grammar.\\
\bottomrule
\end{tabular}
\caption{Frontmatter separation of current release status, OP5 local closure, and historical theorem memory.}
\label{tab:frontmatter-human-llm-architecture}
\end{table}


\subsection*{DOI discipline for the human-facing edition}
\addcontentsline{toc}{subsection}{DOI discipline for the human-facing edition}
\label{subsec:doi-discipline-human-edition}

\begin{center}
\small
\begin{tabular}{>{\raggedright\arraybackslash}p{0.26\textwidth} >{\raggedright\arraybackslash}p{0.28\textwidth} >{\raggedright\arraybackslash}p{0.34\textwidth}}
\toprule
\textbf{Record} & \textbf{Status} & \textbf{Reader-facing rule}\\
\midrule
Concept DOI & stable Companion-line pointer & \href{https://doi.org/10.5281/zenodo.19210728}{10.5281/zenodo.19210728}\\
Version DOI & current public version-specific release DOI & \href{https://doi.org/10.5281/zenodo.20117361}{10.5281/zenodo.20117361}.\\
Prior release metadata & provenance layer only & Historical DOI records remain reconstructible in audit and release files, not in the proof spine.\\
\bottomrule
\end{tabular}
\end{center}

\noindent\textbf{DOI and versioning rule.} DOI records and release identifiers are metadata. They must not be used as mathematical indices, route tags, theorem labels, tensor subscripts, obstruction names, or proof data in the reader-facing argument.

\clearpage
